\begin{proof}[Proof of \cref{obj:lem:setwise-domination-unit-bound}]
Fix \(q>1\).  For \(m\in\mathbb N\), put
\[
A_m(q)=\{x:q<f(x)\le m\}.
\]
This set is measurable, and \(f\mathbf 1_{A_m(q)}\) is integrable because \(|f|\le m\) on \(A_m(q)\) and \(\mu\) is finite.  On this slice,
\[
q\,\mu(A_m(q))
=
\int_{A_m(q)}q\,d\mu
\le
\int_{A_m(q)}f\,d\mu
\le
\left|\int_{A_m(q)}f\,d\mu\right|
\le
\mu(A_m(q)).
\]
Since \(q>1\), this gives \(\mu(A_m(q))=0\) for every \(m\). % lean: abs_le_one_ae_of_setIntegral_le
If \(f(x)>q\), choose \(m\in\mathbb N\) with \(f(x)\le m\); then \(x\in A_m(q)\).  Hence
\[
\mu\{x:f(x)>q\}=0,
\qquad\text{so}\qquad
f\le q\quad \mu\text{-a.e.}
\]
% lean: abs_le_one_ae_of_setIntegral_le

The lower tail is identical after reversing signs.  For \(m\in\mathbb N\), set
\[
B_m(q)=\{x:-m\le f(x)<-q\}.
\]
Again \(f\mathbf 1_{B_m(q)}\) is integrable, and
\[
\int_{B_m(q)}f\,d\mu
\le
\int_{B_m(q)}(-q)\,d\mu
=
-q\,\mu(B_m(q)).
\]
The assumed domination also gives
\[
-\mu(B_m(q))
\le
\int_{B_m(q)}f\,d\mu.
\]
Therefore \((q-1)\mu(B_m(q))\le0\), and \(\mu(B_m(q))=0\).  Taking \(m\) with \(-f(x)\le m\) shows
\[
\mu\{x:f(x)<-q\}=0,
\qquad\text{so}\qquad
-q\le f\quad \mu\text{-a.e.}
\]
% lean: abs_le_one_ae_of_setIntegral_le

Apply the preceding two conclusions to
\[
q_j=1+\frac{1}{j+1},\qquad j=0,1,2,\ldots .
\]
Outside a null set, \(-q_j\le f(x)\le q_j\) for every \(j\).  Letting \(j\to\infty\) gives
\[
-1\le f(x)\le 1
\]
outside that null set, equivalently \(|f(x)|\le1\) \(\mu\)-almost surely. % lean: abs_le_one_ae_of_setIntegral_le
\end{proof}
